% Native author continuation, 2026-10-20.
% Dependencies: 260, 270, 275, 279--283, 285--286, 316--317, 319.
% The designated zero-clearance mark is not a collision.
\section{Exponential payment of positive physical critical fans at unvisited marks}
\label{sec:pf-positive-fans}

The regular boundary estimate does not include the incoming fan at a
critical endpoint action.  We now pay the complete positive fan at an
original zero-clearance mark when all limiting \emph{visited} incidences
are positive.  In particular the fixed-clearance angular slice need
not contain a regular interval at the critical state.  Simultaneous
unvisited tangencies and original section cuts are retained.  The
mechanism is a transfer of the physical branch-supremum estimate from
nearby positive clearance levels; it is not an evaluation of an
almost-everywhere kernel at a singular source point.

Fix $R\in[R_-,R_+]$, $m\ge1$ and $0\le j<m$.  A designated
critical marked germ consists of the original $m$-flight broken path,
its strict visited word, and one nonincident disk tangent to flight
$j$.  All its visited incidences are positive and $p_0=p_m=0$.
The tangency at the mark is approached from $u>0$ and adds no
collision.  Other unvisited tangencies may occur on any of the
flights.  At such a limiting state no derivative of a physically
defined singular collision map is asserted: the smooth functions
below are the one-sided auxiliary action and optical functions of
Corollary~\ref{cor:cr-one-sided}.

Translate the marked center to zero, and denote its oriented line
coordinates by $(\vartheta,\varsigma,u)$, with
$\varsigma\in\{-1,1\}$ and $u=0$ at the germ.  The translation is
only the one-copy parametrization of Theorem~\ref{thm:cf-flux}.
The physical initial state is recovered by following the line through
$j+1$ actual backward collisions and then translating its initial
disk to the original cell.

\subsection{Actual positive sectors and a common nearby clearance level}

For each germ $x$, let $E_x$ be its physical side with the original
initial and terminal section tests, exact occupation and displacement
record, and any fixed original roof interval retained.  All
unvisited-obstacle inequalities are part of $E_x$.  Unite all sector
components and all coordinate copies that represent the same marked
germ before assigning a coefficient.  At a simultaneous cut this
union can have positive opening although its intersection with
$u=0$ is just the limiting point.

\begin{lemma}[Finite analytic sign data on the true physical side]
\label{lem:pf-sign-data}
At a critical marked germ with positive visited incidences there is
a real-analytic inverse chart in $z=(p_0,p_m)$ near $z=0$ for its
fixed strict visited word.  Its actual physical subset $E_x$ is a
finite Boolean combination of real-analytic inequalities.  The
indicator $\one_{E_x}(r e_\theta)$ consequently has, outside a
finite set of directions, a limit as $r\downarrow0$.  The set
$S_x^z$ on which this limit is one is a finite union of angular
intervals, up to their endpoints.  On every compact subinterval of
these intervals the entire sufficiently short punctured radial cone
consists of actual regular trajectories satisfying the original
record tests.  This is an assertion about the physical and original
finite record cuts, not a ray-persistence assertion for an arbitrary
Borel pairing selector.
\end{lemma}
\begin{proof}
On the designated word each disk intersection is transverse.
The circle intersection formulas and specular reflection are
therefore analytic in their local incoming state and parameters.
The two first hits are the prescribed strictly met disks on opposite
sides of the marked foot; the marked disk is not reflected there.
Iterating the prescribed regular intersections gives an analytic
auxiliary map in $(\vartheta,u)$ on a neighborhood of the limiting
point.  All positive visited incidences stay positive on a smaller
neighborhood.  On the actual side this map is exactly the original
collision record, including any later genuine visit to the marked
disk.  By \eqref{eq:bf-positive-determinant} its map to $z$ has
nonzero determinant at the germ.  The analytic inverse-function
theorem gives the asserted inverse chart.  The signed optical
identities extend to this positive-incidence limit by continuity.

Only finitely many disks need be tested near this fixed finite
word.  Indeed its finitely many auxiliary segments stay in one
fixed bounded region after shrinking the neighborhood; disks
outside a sufficiently large bounded region cannot meet a segment.
This uses every candidate center, rather than a sampled subset.
For a nonincident disk whose limiting clearance is zero, the foot
is strictly interior and separated from the incident endpoints by
Lemma~\ref{lem:cf-foot}.  Its signed impact has magnitude $R>0$,
so its clearance is analytic near the germ.  Requiring it to be
positive is precisely the local no-extra-collision condition.
All other nonincident disks are separated by a positive margin
and remain so on a smaller neighborhood.  Thus the simultaneous
physical restrictions are finitely many analytic inequalities; no
forbidden part of the auxiliary continuation is admitted.

At fixed $R$ the original section is a finite union of coordinate
rectangles.  Each of its finitely many tests through collision $m$
is an analytic coordinate inequality on this regular auxiliary
word.  The exact occupation count is a finite Boolean combination
of those tests, with the initial visit counted and terminal visit
tested separately.  The lifted displacement is constant on the
word.  The physical roof is the analytic sum of the actual flight
lengths, so a fixed roof interval adds two analytic inequalities.
No different clock or conditional event is used.  Local angular
lifts remove coordinate seams without changing these tests.

Here is the sign argument, including possible degeneracies.  If
an analytic defining function $a$ is not identically zero, let
$P_d$ be its first nonzero homogeneous Taylor term.  Uniformly on
the circle,
\[
                  a(r e)=r^d\{P_d(e)+O(r)\}.
\]
The restriction of the nonzero homogeneous polynomial $P_d$ to the
circle has finitely many zeros: it is a nonzero trigonometric
polynomial.  Its sign controls the sign of $a(r e)$ away from
those directions, uniformly on compact subsets.  An identically
zero function instead has the constant truth value specified by
its actual inequality convention.  Take the union of the finitely
many exceptional direction sets and apply the original Boolean
formula.  This proves the asserted angular convergence and the
compact cone property.  The genuine active clearance and section
functions are not identically zero: for a clearance this follows
also from the nonzero line-coordinate determinant, and a section
coordinate is pulled back by a locally invertible regular collision
map.  Along an admitted generic ray these physical inequalities
are strict.  Any degenerate additional roof inequality is governed
by its same first nonzero Taylor term.  Thus admitted ray points
are interior points of the actual fixed-word source and of the
nonconstant record tests.
\end{proof}

Call a germ \emph{positive} when $S_x^z$ has positive angular
measure.  The corresponding normalized-action sector in the
$y=G_x^{1/2}(q-q_x)$ coordinates has positive opening exactly in
this case: the invertible linear angular map in
Lemma~\ref{lem:mc-angle} preserves angular null sets.  A zero-opening
germ has zero leading physical fan coefficient by dominated
convergence in Theorem~\ref{thm:cr-critical-density}; this statement
uses the bounded positive-incidence coarea coefficient, not merely
the collision-area nullness of the limiting point.

\begin{lemma}[Common positive-clearance representatives]
\label{lem:pf-representatives}
For any finite collection of positive marked germs, after their
one-copy identifications, and every sufficiently small common
$u>0$, one can choose one actual regular representative $x(u)$
of each germ with marked clearance exactly $u$, satisfying its
original finite record tests, and with $x(u)\longrightarrow x$.
The representatives can be chosen in the interiors of both
one-sided partitions in Theorem~\ref{thm:mb-branch-sup}.
All the positive optical scalars $A_\pm^p$, $h_\pm$ and $u_x$
at these representatives converge to their auxiliary values at
the respective limiting germs.
\end{lemma}
\begin{proof}
In the inverse chart of Lemma~\ref{lem:pf-sign-data}, the marked
clearance $u=u(z)$ is analytic, $u(0)=0$, and $D u(0)\ne0$:
it is a coordinate of the inverse of the nonsingular map
$(\vartheta,u)\mapsto(p_0,p_m)$.  Choose an interior direction
$e$ of a positive physical sector, avoiding all exceptional
directions in that lemma and the two zeros of $D u(0)e$.
Along its admitted physical ray $u(r e)>0$ for small $r>0$;
therefore $d_x:=D u(0)e>0$.  Taylor expansion gives
\[
       u(r e)=d_xr+O(r^2),\qquad
       \frac{d}{dr}u(r e)=d_x+O(r)>0
\]
on a sufficiently short interval.  Invert this strictly increasing
function.  It produces a genuine representative at every
$0<u<u_x^*$, not just along a sequence of levels.  All original
tests hold by the compact cone conclusion of the previous lemma.

Every such point has a nonempty open angular neighborhood on its
fixed-$u$ slice lying in the same actual record source: all its
visited incidences and unvisited clearances are positive and the
nonconstant original analytic tests are strict.  The union of the
exceptional angular sets of the forward and backward partitions
in Theorem~\ref{thm:mb-branch-sup} is null.  Perturb the angular
coordinate within that open slice, by less than $u$ if needed,
to avoid both exceptional sets.  The perturbed representative
still tends to the same germ and is in the interiors of both
actual partition intervals.  No Borel pairing selector is assumed
to persist under this perturbation; such selectors are bounded
later by the original physical fan.

For finitely many germs take the minimum of their positive
$u_x^*$ and of the finitely many required chart sizes.  This proves
the common-level assertion.  Each limiting visited incidence is
positive, so every finite kick and flight factor in
\eqref{eq:bf-half-product} is continuous and finite on its fixed
word up to this limit.  Its entries, the strictly positive
$A_\pm^p$, the ratios $h_\pm$ and the positive determinant formula
\eqref{eq:bf-root-weight} consequently converge as stated.
This local construction is made at fixed $R,m$; no uniform
continuation radius is claimed or required.
\end{proof}

\subsection{One-copy limiting roots and the half-source capacity}

\begin{lemma}[Distinct positive germs use distinct half-source intervals]
\label{lem:pf-distinct}
Distinct one-copy positive germs at fixed $R,m,j$ have distinct
limiting marked parameters $(\varsigma,\vartheta)$.
For a finite family and sufficiently small common $u>0$, the
representatives of Lemma~\ref{lem:pf-representatives} therefore
belong to distinct backward partition intervals, and also to
distinct forward partition intervals, with impact signs as labels.
\end{lemma}
\begin{proof}
Fix the translated marked center, its sign and direction at
zero clearance.  Continue straight from its foot until the first
\emph{strictly} met disk in each direction.  At the limiting
configuration any earlier contact is an unvisited tangency; it
cannot be an alternative positive-incidence visited collision.
The first strict disk and its intersection point are unique.
Specular reflection there fixes the next ray, and the same
argument iterates through the prescribed finite positive numbers
of strict visited collisions.  Thus the two half-paths and their
initial normalization are unique.  Any proposed alternative
strict word using a tangent contact would have a visited zero
incidence and is outside the stated class.  This is the same
one-copy limiting uniqueness used in
Theorem~\ref{thm:mc-critical-sectors}, now based at the marked foot.
It is a uniqueness assertion for limits of actual words, not a
definition of a physical map through grazing.  Multiple physical
sector components at this one state are consequently united;
they are not different roots.

For the finite family choose one common angular seam avoiding
all its limiting parameters, on both labelled circles.  The
construction and constants of Theorem~\ref{thm:mb-branch-sup}
are uniform in this choice: precisely one seam is used before
either direction is partitioned.  Its exact reversal has the
same seam budget, as in that theorem.  There is now a positive
minimum angular separation between different limiting parameters
of each same-sign subfamily containing at least two elements.
For small $u$ the representatives retain at least half this
separation, whereas their original $p_0$ and $p_m$ both tend to zero.

If two representatives of the same sign lay on one backward
partition interval, its momentum would be strictly monotone and
\eqref{eq:mb-one-sided} would give
\[
 |p_0(x(u))-p_0(\widehat x(u))|
   =\left|\int_{\vartheta(x(u))}^{\vartheta(\widehat x(u))}
                       \partial_\vartheta p_0\,d\vartheta\right|
   \ge k_-a_\circ
           |\vartheta(x(u))-\vartheta(\widehat x(u))|.
\]
The left side tends to zero and the right side stays positive,
a contradiction.  The same calculation with $p_m$ proves the
forward assertion.  Different signs already lie on different
labelled source circles.  No claim that distinct sectors of a
single germ have distinct angular limits is made.
\end{proof}

\begin{theorem}[Complete positive limiting half-weight sum]
\label{thm:pf-half-sum}
On the complete one-copy collection $\mathcal F^+_{m,j,R}$ of
positive marked critical germs just defined, including all
simultaneous original cuts, one has
\begin{equation}\label{eq:pf-half-sum}
 \sum_{x\in\mathcal F^+_{m,j,R}}(A_-^p(x))^{-1}\le H_{\rm nm},
 \qquad
 \sum_{x\in\mathcal F^+_{m,j,R}}(A_+^p(x))^{-1}\le H_{\rm nm}.
\end{equation}
Consequently, with the exact constant $Q_{\rm mb}$ of
\eqref{eq:mb-constants},
\begin{equation}\label{eq:pf-root-sum}
 \sum_{x\in\mathcal F^+_{m,j,R}}u_x
       \le Q_{\rm mb}\,\sigma^{\max(j,m-1-j)},
              \qquad \sigma=47/59.
\end{equation}
Neither estimate requires a nonempty regular interval of the
fixed level $u=0$ at the limiting germ.
\end{theorem}
\begin{proof}
Take any finite subfamily.  At each sufficiently small common
$u>0$, Lemmas~\ref{lem:pf-representatives} and
\ref{lem:pf-distinct} place its representatives in different
actual backward intervals.  Hence
\[
 \sum_x (A_-^p(x(u)))^{-1}
 \le\sum_{I\in\mathcal P^-_{j+1}}
                         \sup_I(A^-_{j+1})^{-1}
 \le H_{\rm nm}.
\]
Apply the identical argument in the forward direction with $m-j$
true collisions.  Continuity of the finite positive optical
products permits $u\downarrow0$ in these finite sums.  Taking
the supremum over finite subfamilies proves
\eqref{eq:pf-half-sum}.  This exhausts the complete positive
source; it does not count words or extend a branch through an
unvisited disk.  The nearby representatives need not share an
endpoint momentum or a physical roof.

The individual optical inequality in \eqref{eq:mb-one-sided}
passes to the same limits and gives
\[
 (A_-^p(x))^{-1}\le E_*\sigma^j,
 \qquad (A_+^p(x))^{-1}\le E_*\sigma^{m-1-j}.
\]
The exact positive determinant \eqref{eq:bf-root-weight} and its
uniform ratio bound yield
$u_x\le C_{\rm foot}/(A_-^p(x)A_+^p(x))$.
Sum using the first inequality of \eqref{eq:pf-half-sum} and
the individual forward bound; the result is
$Q_{\rm mb}\sigma^{m-1-j}$.  Use the opposite direction for
$Q_{\rm mb}\sigma^j$.  Since $0<\sigma<1$, the smaller of
these two bounds is exactly the right side of
\eqref{eq:pf-root-sum}.  Both signs have already been included
in $H_{\rm nm}$.  There is no second sign or collision-word factor.
\end{proof}

\subsection{The original incoming fan and its normalization}

Let $\beta_y(x)$ be the angular fraction of the united physical
sector of a germ in the normalized coordinates of
Theorem~\ref{thm:cr-critical-density}.  Write
\[
 \mathfrak F_{m,j,R}
       =\sum_x\mathfrak J(x)\beta_y(x),\qquad
 \mathfrak J(x)=\frac{\omega(x)}{2c_*\sqrt{\det G_x}}.
\]
This is the full positive leading incoming-fan charge at the
specified zero-clearance mark.  Original exact records and roof
intervals are retained in the sector or as restrictions of the
critical values.  A Borel fan selector, when its trace is defined,
is a restriction of this same positive measure.  An arbitrary
Borel selector without a trace retains its exact angular function
from Theorem~\ref{thm:cr-critical-density}; no new boundary value
is assigned to it.

\begin{theorem}[Every-depth exponential positive-fan budget]
\label{thm:pf-fan}
For the original family, all $m\ge1$, all $0\le j<m$ and all
original exact-label and physical-roof restrictions,
\begin{equation}\label{eq:pf-fan}
 \mathfrak F_{m,j,R}
   \le C_{\rm pf}\,\sigma^{\max(j,m-1-j)},
 \qquad C_{\rm pf}=\frac{170}{18c_*}\,Q_{\rm mb}.
\end{equation}
The same estimate holds after any further one-copy fan restriction
bounded by one.  It includes every positive-sector germ with
positive limiting visited incidences, even if other unvisited
tangencies or section cuts collapse its $u=0$ angular slice.
Consequently
\begin{equation}\label{eq:pf-all-j}
 \sum_{j=0}^{m-1}\mathfrak F_{m,j,R}
   \le m C_{\rm pf}\sigma^{(m-1)/2},\qquad
 \lim_{m\to\infty}m^2\sup_{R,n,k}
                   \sum_{j=0}^{m-1}\mathfrak F_{m,j,R}=0.
\end{equation}
The full sum includes shallow indices $j=0$ and $j=m-1$; it
does not require a collision-depth cutoff.
\end{theorem}
\begin{proof}
At the germ let $\theta=\arg g$ and
$\chi_x(\theta)=e_\theta^{\perp\mathsf T}G_xe_\theta^\perp$.
The incoming Gauss-angle fan has its original positive density
\[
              \frac{\omega(x)}{4\pi c_*\chi_x(\theta)}
                           \,d\theta
\]
on its actual sector, before any bounded source selector is
applied.  The two-dimensional positive-matrix integral gives
\[
 \int_0^{2\pi}\frac{d\theta}{\chi_x(\theta)}
             =\frac{2\pi}{\sqrt{\det G_x}}.
\]
For completeness, rotate to diagonal form with eigenvalues
$\lambda_1,\lambda_2>0$; on each quadrant the substitution
$t=\tan\theta$ reduces the integral to
$\int_0^\infty dt/(\lambda_1+\lambda_2t^2)
 =\pi/(2\sqrt{\lambda_1\lambda_2})$.
Thus the full-circle coefficient is exactly $\mathfrak J(x)$;
the physical sector gives $\mathfrak J(x)\beta_y(x)$ by the
actual angular change of variables in Lemma~\ref{lem:mc-angle}.
All components at the same marked state are united before this
bound, so their total is at most one full circle.

The exact inverse-weight identity and the critical Hessian bound
in \eqref{eq:mc-weight}--\eqref{eq:mc-Hbounds} give
\[
 \mathfrak J(x)
   =\frac{\sqrt{\det G_x}}{2c_*R^2}u_x
   \le\frac{170}{18c_*}u_x.
\]
This uses exactly one original conditioning factor.  Its radius
factor is the original $R^{-2}$ in the critical coefficient,
already included in the proved bound
$\sqrt{\det G_x}/R^2\le170/9$; no new probability is normalized.
The equality is the one-sided action coefficient at an unvisited
tangency, not a stability determinant of an undefined singular
return map.

Zero-opening germs have zero leading charge by
Lemma~\ref{lem:pf-sign-data} and the bounded local coarea limit.
For every positive germ discard only its sector fraction on the
upper side, use \eqref{eq:pf-root-sum}, and sum the complete
nonnegative family.  This proves \eqref{eq:pf-fan}.  Bounded
fan selectors can only decrease this positive measure; their
artificial connected components are not counted again.
Finally $\max(j,m-1-j)\ge(m-1)/2$ at every index, so the
full $m$-term sum gives \eqref{eq:pf-all-j}.  Its normalized
upper bound is $C_{\rm pf}m^3\sigma^{(m-1)/2}$, which tends
to zero.  All constants are independent of $R$, the labels and
the roof restrictions, permitting the displayed suprema after
the complete positive sums have been taken.
\end{proof}

\paragraph{Limit order and the exact remaining scope.}
The local operation $u\downarrow0$ in this proof defines and
pays a physical boundary coefficient at a \emph{fixed} radius
and collision count, first for finitely many positive germs.
The resulting bound is uniform, and its nonnegative exhaustion
precedes every collision or radius supremum.  There is no
interchange of a collision limit with a germ-dependent local
chart radius.  In particular no sum of local
$O(\sqrt{s})$ coarea errors on a fixed roof window has been
declared uniformly small.  Equation~\eqref{eq:pf-all-j} concerns
the leading incoming-fan measure itself and its unchanged-roof
pushforward in total variation.  It is not an essential-height
estimate of the entire nearby collision-source density.

The positive-sector hypothesis is verified here for the original
finite physical, section, record and roof cuts by their actual
analytic sign data.  General Borel pairing cuts receive the
stated measure domination when their fan trace is defined, not
an invented analytic description.  Germs with a zero limiting
\emph{visited} incidence are outside the smooth optical-limit
argument; their incoming or pairing traces require their own
proved limit estimate.  Critical fans with no designated
zero-clearance mark, separate section-only or pairing boundary
measures, and the complete shallow same-roof positive-source
shortage also retain their stated obligations.  This theorem
closes the specified full positive marked-fan budget, including
simultaneous unvisited and section corners, without deleting
those sources.  The original raw local limit target
\eqref{eq:v76-target}, physical clock, exact arithmetic factor,
conditional normalization and original order of outer limits
remain unchanged and are not asserted to be completed by this
fan estimate.
